%%%%%%%%%%%%%%%%%%%%%%%%%%%%%%%%%%%%%%%%%
% a0poster Portrait Poster
% LaTeX Template
% Version 1.0 (22/06/13)
%
% The a0poster class was created by:
% Gerlinde Kettl and Matthias Weiser (tex@kettl.de)
%
% This template has been downloaded from:
% http://www.LaTeXTemplates.com
%
% License:
% CC BY-NC-SA 3.0 (http://creativecommons.org/licenses/by-nc-sa/3.0/)
%
%%%%%%%%%%%%%%%%%%%%%%%%%%%%%%%%%%%%%%%%%

%----------------------------------------------------------------------------------------
%	PACKAGES AND OTHER DOCUMENT CONFIGURATIONS
%----------------------------------------------------------------------------------------

\documentclass[a0,portrait]{a0poster}

\usepackage{multicol} % This is so we can have multiple columns of text side-by-side
\columnsep=100pt % This is the amount of white space between the columns in the poster
\columnseprule=3pt % This is the thickness of the black line between the columns in the poster

\usepackage[svgnames]{xcolor} % Specify colors by their 'svgnames', for a full list of all colors available see here: http://www.latextemplates.com/svgnames-colors

\usepackage{times} % Use the times font
%\usepackage{palatino} % Uncomment to use the Palatino font

\usepackage{graphicx,calc} % Required for including images

%\graphicspath{{figures/}} % Location of the graphics files
\graphicspath{{../}} % Location of the graphics files

\usepackage{booktabs} % Top and bottom rules for table
\usepackage[font=small,labelfont=bf]{caption} % Required for specifying captions to tables and figures
\usepackage{amsfonts, amsmath, amsthm, amssymb} % For math fonts, symbols and environments
\usepackage{wrapfig} % Allows wrapping text around tables and figures

\captionsetup[figure]{labelfont={bf},name={Fig.},labelsep=period} % Set picture Fig.1

\makeatletter
\def\@captype{figure}
\makeatother
%\caption{???}\label{xxx} % Use caption and label in \center environment


%wrap figure
\newcommand\myfigure[6]{%
  \ifdim#2>.8\linewidth
    {%
      \centering
      \includegraphics[width=#3]{#4}%
      \captionof{figure}{#5}%
      \label{#6}%
    }%
  \else
  \begin{wrapfigure}{#1}{#2}
    \includegraphics[width=#3]{#4}
    \caption{#5}
    \label{#6}%
  \end{wrapfigure}
  \fi
}


\begin{document}

%----------------------------------------------------------------------------------------
%	POSTER HEADER
%----------------------------------------------------------------------------------------

% The header is divided into two boxes:
% The first is 75% wide and houses the title, subtitle, names, university/organization and contact information
% The second is 25% wide and houses a logo for your university/organization or a photo of you
% The widths of these boxes can be easily edited to accommodate your content as you see fit

\begin{minipage}[c]{0.88\linewidth}
\begin{center}
\Huge \color{MidnightBlue} \textbf{Topological magnons in a one-dimensional itinerant flatband ferromagnet}
\color{Black}\\ % Title
% \Huge\textit{An Exploration of Complexity}\\[2cm] % Subtitle
\Large \textbf{Xiao-Fei Su$^{1,2}$, Zhao-Long Gu$^1$, Zhao-Yang Dong$^1$, Jian-Xin Li$^{1,3*}$}\\[0.5cm] % Author(s)
\Large $^1$Department of Physics, Nanjing University, Nanjing, China 210093 \\[0.4cm]
\Large $^2$School of Physics and Electronic Information, Huaibei Normal University, Huaibei, China 235000\\[0.4cm]
\Large $^3$Collaborative Innovation Center of Advanced Microstructures, Nanjing, China 210093 \\[0.4cm] % University/organization

%\Large \texttt{john@LaTeXTemplates.com} --- 1 (000) 111 1111\\
\end{center}
\end{minipage}
%
\begin{minipage}[c]{0.12\linewidth}
\begin{center}
\includegraphics[width=\linewidth]{NJU.pdf}
\end{center}
\end{minipage}

\vspace{1cm} % A bit of extra whitespace between the header and poster content


%----------------------------------------------------------------------------------------
%	ABSTRACT
%----------------------------------------------------------------------------------------
\begin{minipage}[c]{\linewidth}
\color{MidnightBlue}
\begin{huge}Abstract\end{huge}
\large
\ \ Different from previous scenarios that topological magnons emerge in local spin models, we propose an alternative that itinerant electron magnets can host topological magnons. A one-dimensional Tasaki model with a flatband is considered as the prototype. This model can be viewed as a quarter-filled periodic Anderson model with impurities located in between and hybridizing with the nearest-neighbor conducting electrons, together with a Hubbard repulsion for these electrons. By increasing the Hubbard interaction, the gap between the acoustic and optical magnons closes and reopens while the Berry phase of the acoustic band changes from 0 to $\pi$, leading to the occurrence of a topological transition. After this transition, there always exist in-gap edge magnonic modes, which is consistent with the bulk-edge correspondence. The Hubbard interaction-driven transition reveals a new mechanism to realize nontrivial magnon bands.
\end{minipage}


\vspace{1cm} % A bit of extra whitespace between the header and poster content

%----------------------------------------------------------------------------------------

\begin{multicols}{2} % This is how many columns your poster will be broken into, a portrait poster is generally split into 2 columns


%----------------------------------------------------------------------------------------
%  Introduction
%----------------------------------------------------------------------------------------
\color{SaddleBrown} % SaddleBrown color for the introduction
\section*{Introduction}
\par Band structure with nontrivial topology has been one of the most active fields in condensed matter physics since the discovery of topological insulators. Although pioneering works focused on fermionic bands, the underlying concepts, such as Berry phase and Chern number, apply equally to systems that host bosonic ones. Magnons, the bosonic quanta of collective spin-1 excitations in a magnetically ordered system, also exhibit band structures in crystals. Recently, the search of topological magnons has attracted much attention \cite{OIKSNT_S2010,ZRWL_PRB2013,CHFSBNL_PRL2015}. However, all previous theoretical and experimental works are based on local spin models where the linear spin wave theory (LSWT) provides a comprehensive understanding and the Dzyaloshinskii-Moriya interaction is crucial to generate nontrivial magnonic bands. Besides the local spin magnetism, there exists another class of magnetism, the itinerant magnetism, for which the standard LSWT fails. Then, a natural question of fundamental interest follows: Can topological magnons emerge in itinerant magnets? If so, what is the mechanism? Here, we will investigate this with a one-dimensional interacting electronic model with a flatband as the prototype.


%----------------------------------------------------------------------------------------
%  Model
%----------------------------------------------------------------------------------------
\color{DarkSlateGray} % SaddDarkSlateGray color for the model
\section*{Model and method}

\myfigure{R}{0.4\linewidth}{\linewidth}{lattice_poster.pdf}{(a) Schematic illustration of Eq. (\ref{model}). (b) Band structure of the free part of Eq. (\ref{model}) with $\epsilon=\lambda^2/t-2t$. Here, $\lambda=1.4t$. (c) Illustration of the creation of a spin-1 excitation with a center-of-mass momentum $q$ on the ferromagnetic ground state in the PED method.}
{lattice}

\par The prototype is a one-dimensional Tasaki model \cite{T_PTP1998}. As is illustrated in Fig. \ref{lattice}(a), the Hamiltonian reads
\begin{equation}\label{model}
\begin{aligned}
    H   &   =t\sum_{\langle ij\rangle_{AA}}c_{iA}^\dagger c_{jA}+H.c.+ \lambda\sum_{\langle ij\rangle_{AB}}c_{iA}^\dagger c_{jB}+H.c.\\
        &   +\epsilon\sum_i c^\dagger_{iB}c_{iB}+U_s\sum_i n_{iA\uparrow}n_{iA\downarrow}+U_d\sum_in_{iB\uparrow}n_{iB\downarrow}
\end{aligned}
\end{equation}
When $\epsilon=\lambda^2/t-2t$, the free part of the model possesses an exact flat electron band, which is separated from the upper band by a gap equal to $\lambda^2/t$, as shown in Fig. \ref{lattice}(b). Here, $\lambda$ is taken to be larger than both $U_s$ and $U_d$. Then the Hubbard interactions can be projected onto the flatband. As is well known, the ground state of a half-filled flat electron band with Hubbard interactions is an itinerant ferromagnet. Then, the creation of a spin-1 excitation with a center-of-mass momentum $q$ from the fully spin-polarized ground state is just to choose a single particle state with momentum $k_i$, flip the electron's spin on it and move it onto another single particle state with momentum $k_j=k_i-q$, as illustrated in Fig.\ref{lattice}(c). Thus, the total dimension of the Hilbert space of spin-1 excitations with a definite center-of-mass momentum is linear with respect to the system size $N_q$, which is in sharp contrast to the exponential dependence in the usual exact diagonalization method (UED). Therefore, a much larger system can be numerically accessed by this projected exact diagonalization (PED) method.

%----------------------------------------------------------------------------------------
%	Results
%----------------------------------------------------------------------------------------
\color{DarkSlateGray} % DarkSlateGray color for the rest of the content
\section*{Topological itinerant magnons}

\subsection*{Spectra of low-energy spin-1 excitations}

\begin{multicols}{2}

\myfigure{L}{1.0\linewidth}{\linewidth}{spectrum.pdf}{Magnon bands calculated by PED (solid lines) and UED (red dots, $N_q=8$) with (a) $U_s=0.02$, (b) $U_s=0.40$, (c) $U_s$=0.78. Other parameters are fixed at $t=1.0$, $\lambda=1.4$, $U_d=1.0$. The inset in (a) shows the magnon bands calculated by PED with $N_q=800$ around $q=\pi$, where the gap between the acoustic band and optical band is obvious.}{spectrum}

\begin{itemize}
  \item PED results and UED results agree with each other qualitatively, verifying the validity of the PED method.
  \item PED results converge when  $N_q\ge60$ although size dependence is obvious when $N_q<60$ in the case of a small $U_s=0.02$.
  \item There exist two branches of magnon bands, a gapless acoustic branch and a gapped optical branch.
  \item The gap between these two magnon bands closes and reopens with the increase of $U_s$, suggesting a kind of transition.
  \item This transition is not a traditional one because the ground state remains unchanged before and after it.
\end{itemize}
\end{multicols}

\subsection*{Berry phase and phase diagram}

\begin{multicols}{2}

\myfigure{L}{1.0\linewidth}{\linewidth}{gap_berry_and_phase_diagram.pdf}{(a) The blue line is the magnitude of the gap between the acoustic and optical band at $q=\pi$ with $N_q=800$. The green solid/dashed line represents the corresponding Berry phase of the acoustic/optical band. Other parameters are fixed at $t=1.0$, $\lambda=1.4$, $U_d=1.0$. (b) Phase diagram of Eq. (\ref{model}). The dotted line marks the parameters used in Fig.\ref{gap}(a) and in Fig.\ref{edge}(a).}{gap}
\begin{itemize}
  \item With the gap closing and reopening, the Berry phase of the acoustic band changes from $0$ to $\pi$ while that of the optical band changes from $\pi$ to $0$.
  \item The Hubbard interaction induced transition is a topological transition.
  \item $\lambda$ and $U_s$ play a cooperative role in the generation of topological acoustic magnons because the onsite energy of the B sites $\epsilon=\lambda^2/t-2t$ increases with the increase of $\lambda$, which results in transfers of electrons from the B site to A site, and thus enhances the effect of the Hubbard interactions on the A site.
\end{itemize}

\end{multicols}

\subsection*{Edge states}

\begin{multicols}{2}

\myfigure{R}{1.0\linewidth}{\linewidth}{edge.pdf}{At the top: (a) Magnon spectrum of Eq. (\ref{model}) subject to open boundary conditions. At the bottom: Difference of the $S_z$ profile between the in-gap or near-gap magnon states and ground state with (b) $U_s=0.02$, (c) $U_s=0.30$ and (d) $U_s=0.78$. Other parameters are fixed at $t=1.0$, $\lambda=1.4$, $U_d=1.0$.}{edge}

\begin{itemize}
  \item After the topological transition, there always exist in-gap magnonic edge modes, which is consistent with the bulk-edge correspondence.
  \item There seems to be in-gap magnonic edge modes before the topological transition when $U_s\sim0.3$. However, these edge states disappear without closing the gap when $U_s$ decreases to $0.02$. Due to the finite size effect, whether these modes are intrinsic remains unclear.
\end{itemize}

\end{multicols}


\section*{Stability of topological itinerant magnons}

\begin{multicols}{2}
\subsection*{Nonflatness}

\begin{equation}\label{delta}
H_{\Delta}=\Delta\sum_ic_{iB}^\dagger c_{iB}
\end{equation}

\myfigure{R}{1.0\linewidth}{\linewidth}{delta_spectrum.pdf}{(a)-(d) Spin-1 excitation spectra calculated by PED with $N_q=800$ of Eq. (\ref{model}) perturbed by Eq. (\ref{delta}) with (a) $U_s=0.30$, $\Delta=0$, (b) $U_s=0.30$, $\Delta=0.22$, (c) $U_s=0.30$, $\Delta=0.30$, (d) $U_s=0.30$, $\Delta=0.45$. The shaded area represents the Stoner continuum. Insets of (a)-(d) show the corresponding perturbed free electron bands. (e) The phase diagram in the $\Delta$-$U_s$ parameter space. NFM, FM and TFM represent nonferromagnetic phase, ferromagnetic magnons and topological ferromagnetic magnons, respectively. Red stars mark the parameters used in (a)-(d). Other parameters are fixed at $t=1.0$, $\lambda=1.4$, $U_d=1.0$.}{dspectrum}

\subsection*{Nearest-neighbor interactions}

\begin{equation}\label{V}
H_{V}=V\sum_{\langle ij\rangle_{AB}}n_{iA}^\dagger n_{jB}
\end{equation}

\myfigure{R}{1.0\linewidth}{\linewidth}{v_spectrum.pdf}{(a)-(d) Spin-1 excitation spectra calculated by PED with $N_q=800$ of Eq. (\ref{model}) perturbed by Eq. (\ref{V}) with (a) $U_s=0.30$, $V=0$, (b) $U_s=0.30$, $V=0.09$, (c) $U_s=0.30$, $V=0.40$, (d) $U_s=0.30$, $V=0.64$. The shaded area represents the Stoner continuum. (e) The phase diagram in the $V$-$U_s$ parameter space. NFM, FM and TFM represents nonferromagnetic phase, ferromagnetic magnons and topological ferromagnetic magnons, respectively. Red stars mark the parameters used in (a)-(d). Other parameters are fixed at $t=1.0$, $\lambda=1.25$, $U_d=1.0$.}{vspectrum}

\end{multicols}

%----------------------------------------------------------------------------------------
%	CONCLUSIONS
%----------------------------------------------------------------------------------------
\section*{Conclusions}
\par We elaborate a new scenario that topological acoustic magnons can emerge from itinerant flatband ferromagnet, which is different from previous ones that topological magnons exists in local spin models. The prototype model is a one-dimensional flatband ferromagnet, which can be viewed as a quarter filled periodic Anderson model with impurities located at the center of the bonds and hybridizing with conducting electrons at their neighboring sites. Concentrating on the physics related to the flatband, we can project the model Hamiltonian onto the flat band and carry out the large-size calculation of spin-1 excitations based on the exact diagonalization method. We find a correlation-driven topological transition to realize nontrivial acoustic magnons, which is driven by the on-site Hubbard repulsion for conducting electrons.

%----------------------------------------------------------------------------------------
%	REFERENCES
%----------------------------------------------------------------------------------------
\color{DarkSlateGray} % Set the color back to DarkSlateGray for the rest of the content

\nocite{*} % Print all references regardless of whether they were cited in the poster or not

\begin{thebibliography}{99}

\bibitem{OIKSNT_S2010}Y. Onose, T. Ideue, H. Katsura, Y. Shiomi, N. Nagaosa, and Y. Tokura, Science \textbf{329}, 297 (2010).
\bibitem{ZRWL_PRB2013}L. Zhang, J. Ren, J.-S.Wang, and B. Li, Phys. Rev. B \textbf{87}, 144101 (2013).
\bibitem{CHFSBNL_PRL2015}R. Chisnell, J. S. Helton, D. E. Freedman, D. K. Singh, R. I. Bewley, D. G. Nocera, and Y. S. Lee, Phys.Rev.Lett. \textbf{115}, 147201 (2015).
\bibitem{T_PTP1998}H. Tasaki, Prog. Thoer. Phys. \textbf{99}, 489 (1998)

\end{thebibliography}

[*] jxli@nju.edu.cn


%----------------------------------------------------------------------------------------
%	ACKNOWLEDGEMENTS
%----------------------------------------------------------------------------------------
\section*{Acknowledgements}
\par We gratefully acknowledge Zuo-Dong Yu for fruitful discussions. This work was supported by the National Natural Science Foundation of China (11774152) and National Key Projects for Research and Development of China (Grant No. 2016YFA0300401). X.-F. S. and Z.-L. G. contributed equally to this work.

%----------------------------------------------------------------------------------------

\end{multicols}

\end{document}
